\documentclass[11pt]{article}
\usepackage[margin=2.2cm]{geometry}
\usepackage[T1]{fontenc}
\usepackage{amsmath,amssymb,booktabs,graphicx,xcolor,enumitem,caption,float}
\usepackage{hyperref}
\setlist{itemsep=2pt,topsep=3pt}
\hypersetup{colorlinks=true,linkcolor=blue!50!black,urlcolor=blue!50!black}
\captionsetup{font=small,labelfont=bf}
\setlength{\emergencystretch}{3em}
\setlength{\parskip}{4pt}
% every figure sits exactly where its text is ([H]); the text about a figure comes right before it
\newcommand{\fig}[2]{\begin{figure}[H]\centering\includegraphics[width=\textwidth]{figures/#1}\caption{#2}\end{figure}}
\newcommand{\takeaway}[1]{\par\noindent\fcolorbox{gray!50}{gray!8}{\parbox{\dimexpr\textwidth-2\fboxsep-2\fboxrule}{\textbf{Take-away.} #1}}\par\medskip}

\title{Level 1: what the data taught us, and what the next data will test}
\author{ccd\_opt\_ml}
\date{2026-11-10}

\begin{document}
\maketitle

\section*{The story in five steps}
\begin{enumerate}
\item \textbf{First night (v1).} Dark charge and gain are mostly ``which amplifier is it''; a trace names its amplifier, so
  models predicted them by recognising the amplifier. That fails on a new amplifier. Most of the error is one number per
  amplifier, which cancels when two settings of the \emph{same} amplifier are compared: the idea of the change models.
\item \textbf{What moves things (v2).} The integration time \texttt{delayInt} moves both noise and dark charge; after it,
  the vertical clock lows (noise), VSUB and \texttt{delayRGH} (dark charge), VDD and VR (gain).
\item \textbf{What v2 showed.} Change models carry over to new amplifiers; for the noise the traces beat the settings
  (88--90\,\% against 69--72\,\% direction); for dark charge and gain the settings alone already get it.
\item \textbf{The integration time did most of the work}, exactly as the physics says, and in v2 no voltage was ever
  changed without changing the timing.
\item \textbf{The next data (v3)} fixes the timing and moves each voltage on its own, over wider ranges, to measure what
  the voltages do to noise and dark charge (effects down to about 1\,\% and 2\,\%).
\end{enumerate}

\paragraph{Words.} An \emph{image} (about 10 minutes) gives the true \textbf{noise} $\sigma_e$, \textbf{dark charge}
$\lambda$ and \textbf{gain}. A \emph{trace} is a fast (0.2\,s) recording of the output. An \emph{absolute} model predicts
a value; a \emph{change} model predicts how a value changes between two packets of the same amplifier. A \emph{new
amplifier} is one the model never saw in training. ``Direction right'' counts clear changes ($\sigma_e$ $>$ 5\,\%,
$\lambda$ $>$ 15\,\%, gain $>$ 1.2\,\%) whose up/down the model got right.

% ======================================================================================================
\section{What the first night (v1) taught us}

\subsection{Dark charge and gain sort by amplifier}
Each dot is one packet of the first night, per amplifier. The gain sits at a different level for each amplifier, and the
dark charge too; the bar chart on the right says how much of each quantity is explained by knowing only the amplifier
(red) or only the settings (blue). Gain: 0.83 by the amplifier, 0.05 by the settings. Dark charge: 0.31 by the
amplifier, 0.01 by the settings. Noise: about equal (0.25 and 0.27).
\fig{s1_amp_vs_settings.png}{v1: every packet per amplifier, and what explains each quantity.}
\takeaway{Gain and dark charge are mostly a property of the amplifier; the noise depends on the settings as much as on
the amplifier.}

\subsection{A trace names its amplifier}
For every trace feature we measured how much of its spread is just ``which amplifier'' ($\eta^2$, left). 47 of 246
features are amplifier fingerprints ($\eta^2>0.9$), most of them from the reference trace (taken at nominal, so only the
amplifier changes it). A classifier names the amplifier from one trace 98--100\,\% of the time (right). The own-trace
noise features are mostly \emph{not} fingerprints (left, blue at low $\eta^2$): they follow the settings.
\fig{s1_fingerprint.png}{v1: how much of each trace feature is the amplifier's identity, and how well a trace names its amplifier.}
\takeaway{A model given traces always knows which amplifier it is looking at. For gain and dark charge, that is most of
the answer on a known amplifier.}

\subsection{So: good on known amplifiers, poor on a new one, and one image fixes most of it}
A trace model (boosted trees on own and reference trace features) on amplifiers it knows, on a new amplifier, and on a
new amplifier after one image of it (its level measured). Gain: 0.9\,\% on known amplifiers, 4.3\,\% on a new one
(worse than the average, 2.5\,\%), 1.3\,\% after one image. Dark charge: 25.7 / 42.0 / 31.2\,\%. Noise: 10.3 / 14.7 /
10.5\,\%.
\fig{s1_known_new_onefix.png}{v1: typical error of a trace model.}
\takeaway{On a new amplifier the model has learned a lookup with nothing to look up. One image of the new amplifier
restores the noise and the gain; the dark charge only partly.}

\subsection{The way out: predict changes}
Left: the error of the absolute model on a new amplifier, split into the part that is one constant per amplifier (its
offset) and the rest. The offset is most of it (dark charge 36 of 58\,\%, gain 4.2 of 5.1\,\%, noise 12.7 of 20\,\%).
Comparing two settings of the \emph{same} amplifier removes that constant. That is the change model. Right: in v1 every
change was measured against nominal, and nominal is near the best noise, so 96\,\% of the clear noise changes went
\emph{up}: a model could score by always answering ``up''. v2 therefore compares packets with each other (51\,\% up).
\fig{s1_change_idea.png}{Why change models, and why v2 changed how changes are measured.}
\takeaway{Change models were built to cancel the amplifier's level; v2 was designed so that changes go both ways.}

% ======================================================================================================
\section{Which parameters move things (v2)}

\subsection{Correlations: settings against the labels}
Every number is a correlation ($-1$ to $+1$). Left: setting against label over all packets. Right: change of a setting
against change of the label, for two packets of one amplifier (what the change models learn). The strongest by far is
\texttt{delayInt}: noise $-0.68$, dark charge $+0.89$ (changes). Next: for the noise the vertical clock lows
(\texttt{V1\_B\_L}, which in v2 moved together with the other V lows: $+0.42$) and \texttt{SW\_L} ($-0.29$); for the gain
VDD ($-0.61$) and VR ($-0.55$); for the dark charge \texttt{RG\_L} ($-0.27$), \texttt{delay\_H} ($+0.19$) and VSUB.
\fig{s2_corr_settings.png}{v2: correlation of each setting with each label (levels left, changes right).}

\subsection{Correlations: settings against what the trace measures}
The same, against key trace features (left of the black line) and the labels (right). \texttt{delayInt} drives the
trace's noise power in every band ($+0.9$) and its white noise per skip ($-0.5$); \texttt{delayRGH} the
pedestal-window settling ($+0.9$); \texttt{delay\_H} the baseline droop ($-0.7$); VR the reset-level scatter ($-0.5$).
The vertical clock lows, which move the image noise ($+0.4$), barely show in any of these trace features.
\fig{s2_corr_traces.png}{v2: correlation of each setting with trace features and labels.}

\subsection{The effect of each setting, measured and as the best change models see it}
All settings fitted together: the change of the label for a one-standard-deviation change of each setting, from the
images (black) and from the predictions of two change models on a new amplifier (blue: settings only; green: traces
only). Noise: \texttt{delayInt} $-13\,\%$, the vertical clock lows $+5.2\,\%$, VSUB $+2.4\,\%$. The trace model sees the
clock-low effect ($+3.1\,\%$), the settings model does not ($-0.3\,\%$). Dark charge: \texttt{delayInt} $+27\,\%$, VSUB
$+7.4\,\%$, \texttt{delayRGH} $+5.9\,\%$; the settings model sees VSUB ($+6.8\,\%$), the trace model hardly
($+1.6\,\%$).
\fig{s2_model_effects.png}{v2: per-setting effect on a change, measured and predicted.}
\takeaway{\texttt{delayInt} first, by far. Second for the noise: the vertical clock low levels (which the traces pick
up); for the dark charge: VSUB and \texttt{delayRGH} (which only the settings carry); for the gain: VDD and VR.}

\subsection{What matters most: the top five in every view}
For each label, the five settings that matter most in each view: the measured effect and the correlations (from the
images), and what each model relies on or predicts (new amplifier). A setting that tops every row is a robust result;
one that appears in only one row is that model's or that view's peculiarity. The last table lists the trace features
the trace models use most. How to read the rows:
\begin{itemize}
\item \textbf{measured effect}: change of the label for a one-standard-deviation change of the setting, all settings
  fitted together (so a setting gets credit only for what the others do not explain);
\item \textbf{correlations}: how strongly the setting and the label move together, one setting at a time ($-1$ to $+1$);
\item \textbf{relevance}: how much the Gaussian process uses the setting (1 / length scale, the largest set to 1);
\item \textbf{predicted effect}: the same fit as the measured effect, on the model's predictions instead of the images:
  what the model believes each setting does.
\end{itemize}
For the noise, \texttt{delayInt} is first in every row and the vertical clock lows second in four of seven; for the dark
charge, \texttt{delayInt} first everywhere, then mostly VSUB and \texttt{delayRGH} (the correlations put RG\_L second);
for the gain, VR and VDD first and second in every row except the settings GP, which did not learn the gain at all
(without the amplifier's name, the settings explain almost none of it).
\begin{table}[H]\centering\small
\caption{What matters most for the noise $\sigma_e$: the five largest settings in each view (v2; changes = two packets of one amplifier; new amplifier for the models). V lows / H lows: the vertical / horizontal clock lows, which moved together in v2; 1 SD of a setting = its standard deviation over the v2 packets.}
\begin{tabular}{p{5.2cm}lllll}\toprule
view & 1st & 2nd & 3rd & 4th & 5th\\\midrule
Measured effect on a change (\%/SD, all settings fitted together) & delayInt $-$13.1 & V lows +5.2 & VSUB +2.4 & RG\_H +2.2 & SW\_L $-$1.9\\
Correlation: setting change vs label change & delayInt $-$0.68 & V lows +0.42 & SW\_L $-$0.29 & VSUB +0.20 & TG $-$0.14\\
Correlation: setting vs label (levels) & delayInt $-$0.59 & V lows +0.38 & SW\_L $-$0.36 & H lows $-$0.16 & RG\_L +0.10\\
Settings GP (absolute): relevance (largest = 1) & delayInt 1.00 & V lows 0.19 & TG 0.14 & H lows 0.13 & delayRGH 0.08\\
Settings change GP: relevance (largest = 1) & delayInt 1.00 & VDD 0.52 & TG 0.25 & RG\_L 0.23 & RG\_H 0.22\\
Settings change GP: predicted effect (\%/SD) & delayInt $-$12.4 & delay\_H $-$2.6 & RG\_L +1.6 & H lows $-$1.6 & VDD +1.3\\
Trace combined trees: predicted effect (\%/SD) & delayInt $-$9.4 & SW\_L $-$4.0 & V lows +3.1 & VR $-$2.7 & RG\_L $-$1.5\\
\bottomrule\end{tabular}\end{table}

\begin{table}[H]\centering\small
\caption{What matters most for the dark charge $\lambda$: the five largest settings in each view (v2; changes = two packets of one amplifier; new amplifier for the models). V lows / H lows: the vertical / horizontal clock lows, which moved together in v2; 1 SD of a setting = its standard deviation over the v2 packets.}
\begin{tabular}{p{5.2cm}lllll}\toprule
view & 1st & 2nd & 3rd & 4th & 5th\\\midrule
Measured effect on a change (\%/SD, all settings fitted together) & delayInt +26.7 & VSUB +7.4 & delayRGH +5.9 & delay\_H +3.6 & VDD +2.1\\
Correlation: setting change vs label change & delayInt +0.89 & RG\_L $-$0.27 & delay\_H +0.19 & TG +0.18 & H lows +0.13\\
Correlation: setting vs label (levels) & delayInt +0.63 & RG\_L $-$0.20 & H lows +0.17 & VSUB +0.16 & delay\_H +0.07\\
Settings GP (absolute): relevance (largest = 1) & delayInt 1.00 & VSUB 0.21 & delayRGH 0.20 & V lows 0.01 & H lows 0.01\\
Settings change GP: relevance (largest = 1) & delayInt 1.00 & delayRGH 0.31 & VSUB 0.29 & delay\_H 0.22 & VR 0.14\\
Settings change GP: predicted effect (\%/SD) & delayInt +26.9 & VSUB +6.8 & delay\_H +3.7 & delayRGH +3.7 & VR $-$2.6\\
Trace combined trees: predicted effect (\%/SD) & delayInt +25.5 & delayRGH +4.3 & RG\_L $-$2.4 & TG $-$2.1 & delay\_H +2.1\\
\bottomrule\end{tabular}\end{table}

\begin{table}[H]\centering\small
\caption{What matters most for the gain: the five largest settings in each view (v2; changes = two packets of one amplifier; new amplifier for the models). V lows / H lows: the vertical / horizontal clock lows, which moved together in v2; 1 SD of a setting = its standard deviation over the v2 packets.}
\begin{tabular}{p{5.2cm}lllll}\toprule
view & 1st & 2nd & 3rd & 4th & 5th\\\midrule
Measured effect on a change (\%/SD, all settings fitted together) & VR $-$0.7 & VDD $-$0.6 & delay\_H +0.2 & delayInt $-$0.2 & delayRGH $-$0.2\\
Correlation: setting change vs label change & VDD $-$0.61 & VR $-$0.55 & delayRGH $-$0.21 & delay\_H +0.12 & delayInt $-$0.09\\
Correlation: setting vs label (levels) & VDD $-$0.30 & VR $-$0.28 & delayRGH $-$0.08 & delay\_H +0.07 & RG\_H +0.04\\
Settings GP (absolute): relevance (largest = 1) & \multicolumn{5}{l}{\textit{not learned: the GP fitted this label as noise (all length scales at their start value)}}\\
Settings change GP: relevance (largest = 1) & VR 1.00 & VDD 0.68 & delay\_H 0.37 & delayRGH 0.28 & delayInt 0.28\\
Settings change GP: predicted effect (\%/SD) & VR $-$0.6 & VDD $-$0.6 & delay\_H +0.2 & delayInt $-$0.2 & delayRGH $-$0.1\\
Trace combined trees: predicted effect (\%/SD) & VR $-$0.5 & VDD $-$0.4 & RG\_H +0.2 & delayRGH $-$0.1 & V lows +0.1\\
\bottomrule\end{tabular}\end{table}

\begin{table}[H]\centering\small
\caption{The trace features the trace models use most (permutation importance; ``ref.:'' = the reference trace taken at nominal).}
\begin{tabular}{p{3.6cm}p{2.3cm}p{2.3cm}p{2.3cm}p{2.3cm}p{2.3cm}}\toprule
model, label & 1st & 2nd & 3rd & 4th & 5th\\\midrule
trace trees (absolute), noise $\sigma_e$ & signal-window settling (ADU) & line frequency, whole trace (k\dots{} & noise of 10-skip average & noise power 0.1-1 MHz & row-start recovery time (skips)\\
trace trees (absolute), dark charge $\lambda$ & PC1 variance fraction & ref.: cycle variation in signal wind\dots{} & drift (ADU/ms) & in-window step rate (\%) & waveform modes for 90 \% of the\dots{}\\
trace trees (absolute), gain & ref.: cycle variation in signal wind\dots{} & dip per volt SW (ADU/V) & cycle variation in signal wind\dots{} & ref.: ringing rms (ADU) & dip recovery time (us)\\
trace combined trees (change), noise $\sigma_e$ & row-start recovery time (skips) & signal-window settling (ADU) & 1/f corner frequency (kHz) & transient time constant (us) & CDS-slope lever arm (samples)\\
trace combined trees (change), dark charge $\lambda$ & row-start recovery time (skips) & baseline droop time constant (\dots{} & noise power 3-7.5 MHz & ringing frequency (MHz) & recovery amplitude ratio\\
trace combined trees (change), gain & dip recovery time (us) & dip per volt SW (ADU/V) & Hadamard deviation n=50 (ADU) & baseline coupling (1e-3) & signal-window settling (ADU)\\
\bottomrule\end{tabular}\end{table}



% ======================================================================================================
\section{What v2 showed}

\subsection{Absolute models}
Typical error on amplifiers seen in training (green) and on a new one (red). Noise: the trace physics model gets 14\,\%
on a new amplifier against 20\,\% for the settings and 24\,\% for the average. Dark charge: on a known amplifier the
settings plus the amplifier's name give 5\,\%; on a new amplifier every model is at 44--83\,\%, no better than the average
(49\,\%). Gain: 1--1.6\,\% known, 3--4.5\,\% new.
\fig{s3_absolute.png}{v2 absolute models (session D cross-validation).}

\subsection{One image of the new amplifier}
A new amplifier's error after 0, 1, 2, 4, 8 of its own images. Dark charge: 48 and 53\,\% for the lowest and highest
amplifier (their level is far from the others'), 16\,\% for the middle one; one image brings all three to 17--18\,\%.
\fig{s3_one_image.png}{v2: a new amplifier's error against the number of its images.}

\subsection{Change models}
Direction right on clear changes, new amplifier, random pairs of packets: session D cross-validation (blue) and the
held-out session E (orange). Noise: trace models 85--90\,\% (E: 87--88\,\%) against 69--72\,\% for the settings; the
90\,\% pass mark is met on D and missed by 2 points on E. Dark charge: 96--99\,\% for every model, settings included. Gain:
the settings 94\,\%; the trace models 77--94\,\%.
\fig{s3_change.png}{v2 change models on a new amplifier.}
\takeaway{Change models carry over to new amplifiers. The traces add real information for the noise; for dark charge and
gain the settings are enough.}

% ======================================================================================================
\section{The integration time did most of the work, and we knew it}

\subsection{It does what the physics says}
Noise against \texttt{delayInt}: it falls like white noise averaging down ($\propto 1/\sqrt{\texttt{delayInt}}$, dashed)
and flattens above about 1500 cycles, where slow noise takes over. Dark charge: it grows in proportion to the readout
time per pixel (dashed lines, one per amplifier, no fitted shape). Gain: no dependence.
\fig{s4_physics.png}{v2: the labels against delayInt with the physics expectation.}

\subsection{How much it explains}
Share of the spread explained (cross-validated). The amplifier alone explains 0.19 of the noise, 0.36 of the dark charge
and 0.70 of the gain. Adding \texttt{delayInt} alone brings noise and dark charge to 0.76 and 0.80; all the other settings
and the traces add nothing beyond that (0.65--0.84). The gain needs VDD and VR (0.84 with all settings).
\fig{s4_shares.png}{v2: share explained by the amplifier, delayInt, all timing, all settings, and traces.}

\subsection{Without it}
Left: direction right against how much \texttt{delayInt} changed. For the noise with small or no \texttt{delayInt}
changes, the settings model is at 53--69\,\%, the trace classifier at 73--98\,\%. Right: a settings model not told
\texttt{delayInt} falls from 76 to 50\,\% (noise) and from 97 to 54\,\% (dark charge) on changes where it moved: a coin
flip.
\fig{s4_without.png}{v2: the change models where delayInt does not change, and without delayInt.}
\takeaway{Most of what the models learned in v2 is the integration time, which physics already predicts. In v2 no packet
changed a voltage with the timing fixed, so v2 cannot tell us what the voltages do on their own.}

% ======================================================================================================
\section{What the next data will test}

\subsection{Which settings moved, and together}
Each square is how strongly two settings moved together over the packets (black = always together). v2 (left): the
three vertical clock lows always together, the three horizontal lows always together, the clock swings never varied.
v3 noise set (right): timing fixed (pink: held at nominal), every voltage moved on its own.
\fig{s5_design_corr.png}{Settings moving together: v2 against the v3 noise set.}

\subsection{Voltage changes at fixed timing, over wider ranges}
Left: packets where a voltage changed and the timing did not: none in v2, 57--84 per voltage in the v3 noise set.
Right: the range covered: VSUB 35--70\,V instead of 40--50\,V; VDD and VR twice as wide, RG and SW three to four times;
the clock swings for the first time.
\fig{s5_coverage.png}{What v3 adds.}

\subsection{How small an effect each design can see}
The smallest effect of each voltage (over its v3 range) that a linear fit of all settings together can tell from zero
(2$\sigma$, 3 amplifiers). v3 noise set: dark charge 1.5--3.2\,\%, noise 0.7--1.4\,\%. v2 could not separate the clock
lows at all (locked together), never varied the swings, and needed 6--11\,\% (dark charge) and 2.4--4.6\,\% (noise) for
VSUB and the readout-stage voltages. TG is held at nominal in the v3 noise set (a TG change is not reset-safe).
\fig{s5_detectable.png}{Smallest detectable voltage effect: v2 design against the v3 noise set.}
\takeaway{v3 measures what each voltage does to noise and dark charge at a fixed integration time, the question v2 could
not answer. The dark-charge add-on (v3d) adds one-voltage-at-a-time sweeps and exposure pairs that separate the dark
current rate from charge made during readout.}

% ======================================================================================================
\clearpage
\appendix
\section{Model-by-model audit}
Each figure is one model, with the checks that suit how that model works. All on v2 session D, new amplifier.

\subsection{Settings Gaussian process (absolute)}
A smooth function of the settings, with an error bar. (a--c) predicted against measured with the error bars; (d--f) are
the error bars honest? (misses divided by the stated error bar should follow the black curve); (g--i) which settings it
relies on (1 / length scale). Its error bars are too small on a new amplifier: 45\,\% (noise), 31\,\% (dark charge),
37\,\% (gain) of the measurements lie within $\pm1\sigma$, against 68\,\% if honest; it does not know the amplifier's
level. It relies on \texttt{delayInt}, the V lows and TG (noise), \texttt{delayInt}, VSUB, \texttt{delayRGH} (dark charge).
\fig{appendix_abs_settings_gp.png}{Settings GP (absolute).}

\subsection{Trace physics model (absolute)}
Noise from a formula: $\sigma_e$ = (the trace's white noise per skip $/\sqrt{N}$) $\div$ a learned gain $\times$ a
learned correction; dark charge learned with no formula. (a) The formula's trace term against the image's noise in ADU:
correlation 0.77. (b) The learned gain: 3.0\,\% error. (c) The learned correction: typically 9\,\%. (d) Final noise:
14\,\%. (e) Dark charge: 83\,\% (no formula, amplifier level unknown). (f) Error bars: noise 52/84\,\% within
$\pm1/2\sigma$ (ideal 68/95), dark charge 2/9\,\%: far too confident.
\fig{appendix_abs_physics.png}{Trace physics model (absolute).}

\subsection{Trace boosted trees (absolute)}
(a--c) predicted against measured; (d--f) the trace features it relies on; (g--i) how its prediction moves with its most
important feature. For the dark charge its top features are slow-drift and step features of its own trace and an
edge feature of the reference trace (taken at nominal, so it mostly identifies the amplifier).
\fig{appendix_abs_trees.png}{Trace boosted trees (absolute).}

\subsection{Settings change Gaussian process}
Change of the settings $\to$ change of the label. (a--c) predicted against true change; (d--f) which setting changes it
relies on; (g--i) its learned response to a change of \texttt{delayInt} alone, from three starting values. Noise: 69\,\%
direction, relying on \texttt{delayInt}, VDD, TG; dark charge 97\,\% (\texttt{delayInt}, \texttt{delayRGH}, VSUB); gain
93\,\% (VR, VDD).
\fig{appendix_chg_settings_gp.png}{Settings change GP.}

\subsection{Physics difference model (noise change)}
$\Delta\log\sigma_e = \Delta\log(\text{trace white noise}) - \Delta\log g + \text{correction}$. (a) The trace's white
noise alone: 68\,\% direction; (b) with the learned gain change: 68\,\%; (c) with the learned correction: 86\,\%; (d) the
correction is typically 16\,\%: the white noise alone misses what does not average down.
\fig{appendix_chg_physics.png}{Physics difference model.}

\subsection{Combined trees (change)}
Trees on the features of trace A, trace B and B $-$ A. (a--c) predicted against true change; (d--f) does it get the
size right (slope 0.58 noise, 0.86 dark charge, 0.37 gain: it predicts changes too small); (g--i) which features
(green) and which part (purple): for the dark charge it uses the difference B $-$ A almost entirely; for the noise all
three parts.
\fig{appendix_chg_combined_trees.png}{Combined trees (change).}

\subsection{Direction classifier (change)}
The probability that the label goes up from A to B. (a--c) is its confidence right (on the diagonal = yes; ROC area 0.95
noise, 0.99 dark charge, 0.92 gain); (d--f) its answers for changes that truly went up (red) and down (blue); (g--i)
accuracy when only its confident answers are kept: noise 89\,\% on all, 96\,\% on the most confident half.
\fig{appendix_chg_classifier.png}{Direction classifier (change).}

\section*{Files}
Figures: \texttt{lvl1/story\_figures.py} (sections 1--5), \texttt{lvl1/model\_diagnostics\_v2.py} (appendix), and the
one-image figure from \texttt{lvl1/amplifier\_identity.py}. Numbers: \texttt{story.json},
\texttt{processed\_v2\_trainD/lvl1\_results/model\_diagnostics/}.

\end{document}
